\chapter{多自由度系统与模态分析}
\intro{本章将两自由度的振动理论系统地推广至任意 $n$ 自由度系统。多自由度系统是工程结构动力学分析的基本模型：有限元模型、集中质量模型、连续体离散化模型等最终都归结为 $n$ 自由度系统的运动方程。本章重点讨论模态分析——通过特征值问题将耦合的 $n$ 自由度系统解耦为 $n$ 个独立的单自由度系统，并以此为基础建立模态叠加法、讨论比例阻尼和模态截断技术。}

\section{$n$ 自由度系统的运动方程}

\subsection{矩阵形式的运动方程}

$n$ 自由度线性系统的运动方程与两自由度系统具有完全相同的形式：

\begin{myformula}
\[
\mathbf{M}\ddot{\mathbf{x}} + \mathbf{C}\dot{\mathbf{x}} + \mathbf{K}\mathbf{x} = \mathbf{F}(t)
\]
\end{myformula}

其中各矩阵的分量含义和性质如下。

\subsection{质量矩阵 $\mathbf{M}$ （$n \times n$）}

\begin{itemize}
\item \textbf{集中质量矩阵}：$\mathbf{M} = \text{diag}(m_1, m_2, \ldots, m_n)$，每个对角元对应一个自由度的质量。形式最简单，适用于集中质量模型。
\item \textbf{一致质量矩阵}（有限元法）：$\mathbf{M} = \int \rho \mathbf{N}^{\top}\mathbf{N}\,dV$，非对角，与刚度阵使用相同的形函数 $\mathbf{N}$。
\item \textbf{性质}：对称（$\mathbf{M} = \mathbf{M}^{\top}$），正定（$\mathbf{x}^{\top}\mathbf{M}\mathbf{x} > 0$ 对所有 $\mathbf{x} \neq \mathbf{0}$）。
\end{itemize}

\subsection{刚度矩阵 $\mathbf{K}$ （$n \times n$）}

\begin{itemize}
\item \textbf{构成}：$K_{ij}$ = 在自由度 $j$ 上产生单位位移所需在自由度 $i$ 施加的力（其他自由度位移 = 0）。
\item \textbf{对角元}：$K_{ii} = \sum_{j} k_{ij}$（与自由度 $i$ 连接的所有弹簧刚度之和）。
\item \textbf{非对角元}：$K_{ij} = -k_{ij}$（连接自由度 $i$ 和 $j$ 的弹簧刚度的负值，$i \neq j$），若两自由度间无直接弹性连接则为零。
\item \textbf{性质}：对称（Maxwell 互等定理），半正定（$\mathbf{x}^{\top}\mathbf{K}\mathbf{x} \geq 0$）。当系统完全约束（无刚体自由度）时，$\mathbf{K}$ 为正定;存在刚体模态时 $\mathbf{K}$ 半正定（至少有一个零特征值）。
\item \textbf{带宽}：在有限元模型中，$\mathbf{K}$ 通常是带状稀疏矩阵，带宽取决于节点编号方式。
\end{itemize}

\subsection{阻尼矩阵 $\mathbf{C}$ （$n \times n$）}

阻尼的真实物理机制极为复杂（材料内耗、接头摩擦、空气阻力、辐射阻尼等），工程中通常不直接从第一原理建立 $\mathbf{C}$，而是采用以下简化模型：

\begin{enumerate}
\item \textbf{比例阻尼（Rayleigh 阻尼）}：$\mathbf{C} = \alpha\mathbf{M} + \beta\mathbf{K}$，可与 $\mathbf{M}, \mathbf{K}$ 同时对角化。这是模态叠加法的基础假设。
\item \textbf{模态阻尼}：直接为各阶模态指定阻尼比 $\zeta_i$，不构造物理阻尼矩阵。
\item \textbf{Caughey 阻尼}：$\mathbf{C} = \mathbf{M} \sum_{j=0}^{n-1} a_j (\mathbf{M}^{-1}\mathbf{K})^j$，$\mathbf{C}$ 的任何可被 $\mathbf{\Phi}$ 对角化的充要条件（推广的 Rayleigh 阻尼）。
\item \textbf{非比例阻尼}：不能对角化的 $\mathbf{C}$，需在状态空间中处理或用复模态理论。
\end{enumerate}

\subsection{多自由度弹簧-质量链的 TikZ 示意图}

\begin{tikzpicture}[scale=0.7]
  % 左墙
  \draw[thick] (0,0) -- (0,2.5);
  \fill[gray!30] (-0.5,0) rectangle (0,2.5);
  \foreach \y in {0.2,0.6,...,2.3} {\draw[thick] (0,\y) -- (-0.3,\y+0.15);}
  
  % 多质量-弹簧链
  \foreach \i in {0,1,...,5} {
    % 弹簧
    \pgfmathsetmacro{\x}{1.0 + 2.0*\i}
    \ifnum\i<5
      \draw[thick,decorate,decoration={coil,aspect=0.6,amplitude=3,segment length=6}] 
        (\x+0.5,1.5) -- (\x+1.8,1.5);
      \node at (\x+1.15,2.0) {$k_\i$};
      \node at (\x+1.15,0.7) {$c_\i$};
    \fi
    
    % 质量块
    \pgfmathsetmacro{\mx}{0.5 + 2.0*\i}
    \fill[blue!25] (\mx,0.2) rectangle (\mx+1.0,2.3);
    \draw[thick] (\mx,0.2) rectangle (\mx+1.0,2.3);
    \node at (\mx+0.5,1.25) {$m_\i$};
    \node at (\mx+0.5,-0.3) {$x_\i$};
  }
  
  % 右墙
  \draw[thick] (12.5,0) -- (12.5,2.5);
  \fill[gray!30] (12.5,0) rectangle (13,2.5);
  \foreach \y in {0.2,0.6,...,2.3} {\draw[thick] (12.5,\y) -- (12.8,\y+0.15);}
  
  \node at (6.25,3.2) {$n$ 自由度弹簧-质量链};
\end{tikzpicture}

\section{固有值问题}

\subsection{特征值问题的提法}

无阻尼自由振动（$\mathbf{C} = \mathbf{0}$, $\mathbf{F} = \mathbf{0}$）的方程为 $\mathbf{M}\ddot{\mathbf{x}} + \mathbf{K}\mathbf{x} = \mathbf{0}$。设解为 $\mathbf{x}(t) = \bm{\phi} \sin(\omega t + \theta)$，其中 $\bm{\phi} \in \R^n$ 是振型向量。代入得：

\begin{myformula}
\[
(\mathbf{K} - \omega^2 \mathbf{M}) \bm{\phi} = \mathbf{0}
\]
\end{myformula}

非零解条件 $\det(\mathbf{K} - \omega^2 \mathbf{M}) = 0$ 给出 $n$ 次代数方程（关于 $\omega^2$），可解得 $n$ 个固有频率：

\begin{myformula}
\[
0 \leq \omega_1 \leq \omega_2 \leq \cdots \leq \omega_n
\]
\end{myformula}

等号对应于重频（对称系统的退化模态）。

将 $\omega_i^2$ 代回特征方程求相应的振型向量 $\bm{\phi}_i$（确定至比例因子）。所有 $n$ 个振型构成 \textbf{模态矩阵}：

\begin{myformula}
\[
\mathbf{\Phi} = [\bm{\phi}_1, \bm{\phi}_2, \ldots, \bm{\phi}_n] \in \R^{n\times n}
\]
\end{myformula}

$\mathbf{\Phi}$ 的第 $i$ 列是第 $i$ 阶振型，第 $j$ 行是第 $j$ 个自由度在各阶振型中的分量。

\subsection{固有频率的物理意义}

\begin{itemize}
\item $\omega_1$：基频（fundamental frequency），全系统的最低固有频率。对应振型一般是所有自由度同向运动。
\item 随着阶次增加，振型中出现越来越多的"节点"（零位移点），相邻质量间的相位反转增多。
\item 第 $i$ 阶振型有 $i-1$ 个内部节点（两端不算）。
\item $\omega_n$：最高固有频率，对应最"曲折"的振型——相邻自由度反向剧烈运动。
\end{itemize}

\subsection{前三阶振型的 TikZ 图}

\begin{tikzpicture}[scale=0.8]
  % 第一阶
  \node at (5,5.5) {\textbf{第一阶振型} ($\omega_1$，同向运动)};
  \draw[thick] (0,5.0) -- (10,5.0);
  \foreach \i in {0,1,...,4} {
    \fill[blue!40] (0.5+2*\i,4.7) rectangle (1.5+2*\i,5.3);
    \node at (1+2*\i,5.0) {$m_\i$};
    \draw[-Stealth,blue,thick] (1+2*\i,5.5) -- (1+2*\i,5.9);
  }
  \node at (5,4.3) {全部同相};
  
  % 第二阶
  \node at (5,3.8) {\textbf{第二阶振型} ($\omega_2$，一个节点)};
  \draw[thick] (0,3.3) -- (10,3.3);
  \foreach \i in {0,1,...,4} {
    \fill[blue!40] (0.5+2*\i,3.0) rectangle (1.5+2*\i,3.6);
    \node at (1+2*\i,3.3) {$m_\i$};
  }
  \foreach \i in {0,1,2} {
    \draw[-Stealth,red,thick] (1+2*\i,3.8) -- (1+2*\i,4.2);
  }
  \foreach \i in {3,4} {
    \draw[-Stealth,red,thick] (1+2*\i,3.8) -- (1+2*\i,3.4);
  }
  \node at (3.5,2.7) {左侧同相, 右侧反相};
  \fill[green!60!black] (6.5,3.3) circle (3pt);
  \node at (6.5,2.9) {节点};
  
  % 第三阶
  \node at (5,2.2) {\textbf{第三阶振型} ($\omega_3$，两个节点)};
  \draw[thick] (0,1.7) -- (10,1.7);
  \foreach \i in {0,1,...,4} {
    \fill[blue!40] (0.5+2*\i,1.4) rectangle (1.5+2*\i,2.0);
    \node at (1+2*\i,1.7) {$m_\i$};
  }
  \draw[-Stealth,orange,thick] (1,2.2) -- (1,2.6);
  \draw[-Stealth,orange,thick] (3,2.2) -- (3,1.8);
  \draw[-Stealth,orange,thick] (5,2.2) -- (5,2.6);
  \draw[-Stealth,orange,thick] (7,2.2) -- (7,1.8);
  \draw[-Stealth,orange,thick] (9,2.2) -- (9,2.6);
  \fill[green!60!black] (2.0,1.7) circle (3pt);
  \fill[green!60!black] (6.0,1.7) circle (3pt);
  \node at (2.0,1.3) {节点};
  \node at (6.0,1.3) {节点};
\end{tikzpicture}

\subsection{特征值问题的数值解法}

对于大规模 $n$ 自由度系统（$n = 10^3 \sim 10^7$），直接展开行列式求解是不现实的。实际采用以下数值方法：

\begin{enumerate}
\item \textbf{子空间迭代法（Subspace Iteration）}：同时迭代求前 $r$ 阶特征对。适合大型稀疏矩阵，是最早成功的大规模求解方法之一。
\item \textbf{Lanczos 法}：将 $(\mathbf{K}, \mathbf{M})$ 三对角化后求特征值，收敛极快，是最流行的结构动力学求解方法。
\item \textbf{Jacobi-Davidson 法}：适合内点特征值（在频谱内部找特定频率附近的模态）。
\item \textbf{逆幂迭代（Inverse Power Iteration）}：求最低频或最接近给定频移的特征对。
\end{enumerate}

在 Python/NumPy 中，对于中小规模（$n \leq 10^3$）系统，`scipy.linalg.eigh(K, M)` 使用 LAPACK 的 DSYGVD 驱动（分治法），精度和效率俱佳。

\section{模态正交性}

\subsection{关于 $\mathbf{M}$ 和 $\mathbf{K}$ 的正交性}

模态正交性的推导与两自由度情况完全类似。对于 $i \neq j$：

\begin{myformula}
\[
\bm{\phi}_i^{\top} \mathbf{M} \bm{\phi}_j = 0, \qquad (i \neq j)
\]
\[
\bm{\phi}_i^{\top} \mathbf{K} \bm{\phi}_j = 0, \qquad (i \neq j)
\]
\end{myformula}

对所有 $n$ 个振型写成矩阵形式：

\begin{myformula}
\[
\mathbf{\Phi}^{\top} \mathbf{M} \mathbf{\Phi} = \text{diag}(m_1, m_2, \ldots, m_n) = \mathbf{M}_{\text{modal}}
\]
\[
\mathbf{\Phi}^{\top} \mathbf{K} \mathbf{\Phi} = \text{diag}(k_1, k_2, \ldots, k_n) = \mathbf{K}_{\text{modal}}
\]
\end{myformula}

其中 $m_i = \bm{\phi}_i^{\top}\mathbf{M}\bm{\phi}_i$ 为第 $i$ 阶模态质量，$k_i = \bm{\phi}_i^{\top}\mathbf{K}\bm{\phi}_i$ 为第 $i$ 阶模态刚度。模态质量与模态刚度之间的关系：

\[
\omega_i^2 = \frac{k_i}{m_i}
\]

\subsection{质量归一化}

最常见的归一化方式是将所有振型缩放至模态质量为 1：

\[
\tilde{\bm{\phi}}_i = \frac{\bm{\phi}_i}{\sqrt{m_i}}, \qquad m_i = \bm{\phi}_i^{\top}\mathbf{M}\bm{\phi}_i
\]

归一化后的模态矩阵（仍记为 $\mathbf{\Phi}$）满足：

\begin{myformula}
\[
\mathbf{\Phi}^{\top} \mathbf{M} \mathbf{\Phi} = \mathbf{I}_{n\times n}
\]
\[
\mathbf{\Phi}^{\top} \mathbf{K} \mathbf{\Phi} = \text{diag}(\omega_1^2, \omega_2^2, \ldots, \omega_n^2) = \mathbf{\Omega}^2
\]
\end{myformula}

质量归一化的优点在于模态方程中不出现模态质量和模态刚度，固有频率直接出现在方程中。

\subsection{正交性的物理与数学意义}

\begin{itemize}
\item \textbf{物理意义}：不同模态之间没有能量交换。如果以某阶振型激发系统，系统只在该阶振型上振动，不激发其他振型。各阶模态是独立的信息通道。
\item \textbf{数学意义}：正交性使得 $\mathbf{\Phi}$ 成为对 $(\mathbf{K}, \mathbf{M})$ 的 simultaneous diagonalizer（同时对角化矩阵）。这是模态叠加法的数学基础。
\item \textbf{加权内积}：若定义加权内积 $\langle\mathbf{u},\mathbf{v}\rangle_M = \mathbf{u}^{\top}\mathbf{M}\mathbf{v}$，则各振型在该内积下正交。振型构成 $(\R^n, \langle\cdot,\cdot\rangle_M)$ 中的正交基。
\end{itemize}

\section{模态叠加法（Mode Superposition）}

\subsection{坐标变换：物理坐标 $\to$ 模态坐标}

设所有 $n$ 个模态均已求得，做变换：

\begin{myformula}
\[
\mathbf{x}(t) = \mathbf{\Phi} \mathbf{q}(t) = \sum_{i=1}^{n} q_i(t)\,\bm{\phi}_i
\]
\end{myformula}

其中 $\mathbf{q}(t) = [q_1(t), \ldots, q_n(t)]^{\top}$ 为模态坐标（主坐标）。

\subsection{运动方程的解耦}

将 $\mathbf{x} = \mathbf{\Phi}\mathbf{q}$ 代入 $\mathbf{M}\ddot{\mathbf{x}} + \mathbf{C}\dot{\mathbf{x}} + \mathbf{K}\mathbf{x} = \mathbf{F}(t)$，左乘 $\mathbf{\Phi}^{\top}$：

\begin{myformula}
\[
\mathbf{\Phi}^{\top}\mathbf{M}\mathbf{\Phi}\,\ddot{\mathbf{q}} + \mathbf{\Phi}^{\top}\mathbf{C}\mathbf{\Phi}\,\dot{\mathbf{q}} + \mathbf{\Phi}^{\top}\mathbf{K}\mathbf{\Phi}\,\mathbf{q} = \mathbf{\Phi}^{\top}\mathbf{F}(t)
\]
\end{myformula}

若采用质量归一化且 $\mathbf{C}$ 为比例阻尼（可与 $\mathbf{M},\mathbf{K}$ 同时对角化），则 $\mathbf{\Phi}^{\top}\mathbf{C}\mathbf{\Phi} = \text{diag}(2\zeta_1\omega_1, \ldots, 2\zeta_n\omega_n)$，得到 $n$ 个解耦的单自由度方程：

\begin{myformula}
\[
\ddot{q}_i + 2\zeta_i\omega_i\dot{q}_i + \omega_i^2 q_i = \bm{\phi}_i^{\top}\mathbf{F}(t) = f_i(t), \qquad i = 1, 2, \ldots, n
\]
\end{myformula}

其中 $f_i(t) = \bm{\phi}_i^{\top}\mathbf{F}(t)$ 为第 $i$ 阶模态力——外力在第 $i$ 阶振型上的投影。

\subsection{解耦方程的通解}

每个方程都是标准的单自由度系统受迫振动问题。零初始条件下的解（Duhamel 积分）：

\[
q_i(t) = \frac{1}{\omega_{di}} \int_0^t f_i(\tau)\, e^{-\zeta_i\omega_i(t-\tau)} \sin\omega_{di}(t-\tau)\,d\tau
\]

其中 $\omega_{di} = \omega_i\sqrt{1-\zeta_i^2}$。

物理空间的响应为：

\[
\mathbf{x}(t) = \sum_{i=1}^{n} q_i(t)\,\bm{\phi}_i
\]

\subsection{模态叠加法的计算流程}

\begin{center}
\begin{tikzpicture}[scale=0.8, node distance=1.2cm]
  \node[draw,rounded corners,fill=blue!10] (A) {$\mathbf{M}\ddot{\mathbf{x}} + \mathbf{C}\dot{\mathbf{x}} + \mathbf{K}\mathbf{x} = \mathbf{F}(t)$};
  \node[draw,rounded corners,fill=green!10,below of=A,yshift=-0.3cm] (B) {求解特征值问题: $\det(\mathbf{K}-\omega^2\mathbf{M})=0$ 得 $\mathbf{\Phi}$, $\omega_i$};
  \node[draw,rounded corners,fill=yellow!10,below of=B,yshift=-0.3cm] (C) {坐标变换: $\mathbf{x} = \mathbf{\Phi}\mathbf{q}$};
  \node[draw,rounded corners,fill=orange!10,below of=C,yshift=-0.3cm] (D) {解耦: $\ddot{q}_i + 2\zeta_i\omega_i\dot{q}_i + \omega_i^2 q_i = f_i$};
  \node[draw,rounded corners,fill=red!10,below of=D,yshift=-0.3cm] (E) {各模态独立求解 (Duhamel 积分等)};
  \node[draw,rounded corners,fill=blue!10,below of=E,yshift=-0.3cm] (F) {叠加: $\mathbf{x}(t) = \sum_{i=1}^n q_i(t)\,\bm{\phi}_i$};
  
  \draw[-Stealth] (A) -- (B);
  \draw[-Stealth] (B) -- (C);
  \draw[-Stealth] (C) -- (D);
  \draw[-Stealth] (D) -- (E);
  \draw[-Stealth] (E) -- (F);
\end{tikzpicture}
\end{center}

\subsection{坐标变换的 TikZ 示意图}

\begin{tikzpicture}[scale=0.8]
  % 左侧: 物理坐标
  \node at (2,5.0) {\textbf{物理空间 (耦合)}};
  \draw[thick] (0,4.5) -- (4,4.5);
  \fill[blue!30] (0.3,4.2) rectangle (1.0,4.8) node[midway] {\small$m_1$};
  \fill[blue!30] (1.2,4.2) rectangle (1.9,4.8) node[midway] {\small$m_2$};
  \fill[blue!30] (2.1,4.2) rectangle (2.8,4.8) node[midway] {\small$m_3$};
  \fill[blue!30] (3.0,4.2) rectangle (3.7,4.8) node[midway] {\small$m_4$};
  \draw[thick,decorate,decoration={coil,aspect=0.5,amplitude=2,segment length=4}] 
    (1.0,4.5) -- (1.2,4.5);
  \draw[thick,decorate,decoration={coil,aspect=0.5,amplitude=2,segment length=4}] 
    (1.9,4.5) -- (2.1,4.5);
  \draw[thick,decorate,decoration={coil,aspect=0.5,amplitude=2,segment length=4}] 
    (2.8,4.5) -- (3.0,4.5);
  \node at (2,3.9) {弹簧耦合};
  
  % 箭头
  \node at (5.5,4.5) {$\xrightarrow{\mathbf{x} = \mathbf{\Phi}\mathbf{q}}$};
  
  % 右侧: 模态空间
  \node at (9,5.0) {\textbf{模态空间 (解耦)}};
  \node at (7.0,4.5) {$\ddot{q}_1 + \omega_1^2 q_1 = f_1$};
  \node at (7.0,3.6) {$\ddot{q}_2 + \omega_2^2 q_2 = f_2$};
  \node at (11.0,4.5) {$\ddot{q}_3 + \omega_3^2 q_3 = f_3$};
  \node at (11.0,3.6) {$\ddot{q}_4 + \omega_4^2 q_4 = f_4$};
  
  \node at (9.0,3.0) {各模态相互独立};
\end{tikzpicture}

\section{比例阻尼（Rayleigh 阻尼）}

\subsection{为什么需要比例阻尼？}

模态叠加法要求阻尼矩阵 $\mathbf{C}$ 能被 $\mathbf{\Phi}$ 对角化，即：

\[
\mathbf{\Phi}^{\top}\mathbf{C}\mathbf{\Phi} = \text{diag}(\ldots, 2\zeta_i\omega_i, \ldots)
\]

但在工程中我们通常只知道各阶模态的阻尼比 $\zeta_i$（由实验测得或经验估计），而不是物理阻尼矩阵 $\mathbf{C}$。比例阻尼提供了一种从模态阻尼比"反推"物理阻尼矩阵的方法。

\subsection{Rayleigh 阻尼的定义}

\begin{myformula}
\[
\mathbf{C} = \alpha \mathbf{M} + \beta \mathbf{K}
\]
\end{myformula}

其中 $\alpha$（单位：$[\text{s}^{-1}]$）和 $\beta$（单位：$[\text{s}]$）是两个标量参数。

在模态坐标下：
\[
c_i = \bm{\phi}_i^{\top}\mathbf{C}\bm{\phi}_i = \alpha \underbrace{\bm{\phi}_i^{\top}\mathbf{M}\bm{\phi}_i}_{m_i} + \beta \underbrace{\bm{\phi}_i^{\top}\mathbf{K}\bm{\phi}_i}_{k_i} = \alpha m_i + \beta k_i
\]

模态阻尼比：

\begin{myformula}
\[
\zeta_i = \frac{c_i}{2m_i\omega_i} = \frac{\alpha}{2\omega_i} + \frac{\beta\omega_i}{2}
\]
\end{myformula}

\textbf{$\zeta_i$ 与 $\omega_i$ 的关系}：

\begin{itemize}
\item $\alpha$ 项（质量比例阻尼，$\zeta \propto 1/\omega$）：支配低频阻尼。物理上模拟外部粘性阻尼（质量与周围介质的相对运动）。
\item $\beta$ 项（刚度比例阻尼，$\zeta \propto \omega$）：支配高频阻尼。物理上模拟材料内阻尼（应变速率相关的能量耗散）。
\item 两条曲线的交点处阻尼比最小：$\zeta_{\min} = \sqrt{\alpha\beta}$，出现在 $\omega = \sqrt{\alpha/\beta}$。
\end{itemize}

\subsection{参数拟合}

若已知两个模态 $i$ 和 $j$ 的阻尼比 $\zeta_i$, $\zeta_j$，可反解 $\alpha, \beta$：

\[
\begin{bmatrix} \frac{1}{2\omega_i} & \frac{\omega_i}{2} \\[4pt] \frac{1}{2\omega_j} & \frac{\omega_j}{2} \end{bmatrix}
\begin{Bmatrix} \alpha \\ \beta \end{Bmatrix} = 
\begin{Bmatrix} \zeta_i \\ \zeta_j \end{Bmatrix}
\]

通常取第一阶和感兴趣的某一高阶（如对整体响应贡献较大的模态）的两个阻尼比来标定。

\textbf{注意}：Rayleigh 阻尼只有两个参数，只能精确匹配两阶模态的阻尼比，其他模态的阻尼比是推定的（可能偏离实际）。对于需要多阶模态精确阻尼的情况，应使用 Caughey 系列阻尼（$n$ 个参数可精确匹配 $n$ 阶阻尼比）。

\section{模态截断}

\subsection{动机}

对于大规模有限元模型（$n = 10^5 \sim 10^7$），完整模态分析（求全部 $n$ 阶模态）是不可行也不必要的：

\begin{itemize}
\item 计算成本：$O(n^3)$ 求解全部特征值在实际工程规模下无法承受。
\item 物理理由：高频模态的响应贡献极小（惯性力限制+阻尼衰减），在关心的时间尺度上可以忽略。
\item 激励频带限制：若外激励只有低频成分，高频模态根本不被激发。
\end{itemize}

因此只提取前 $r$ 阶模态（$r \ll n$）近似系统响应。

\subsection{截断模态叠加}

用截断的模态矩阵 $\mathbf{\Phi}_r \in \R^{n\times r}$（只含前 $r$ 列）：

\[
\mathbf{x}(t) \approx \sum_{i=1}^{r} q_i(t)\,\bm{\phi}_i = \mathbf{\Phi}_r \mathbf{q}_r(t)
\]

解耦得到 $r$ 个（而非 $n$ 个）单自由度方程。计算复杂度从 $O(n^3)$ 降至 $O(r^3 + nr^2)$。

\textbf{截断准则}：

\begin{enumerate}
\item \textbf{频率准则}：包含所有 $\omega_i \leq \omega_{\max}$ 的模态，其中 $\omega_{\max} \approx 3\sim5 \times \omega_{\text{激励}}$。

\item \textbf{有效模态质量准则}（见下节）：使被截模态的有效质量之和占总质量的足够小比例。

\item \textbf{响应收敛准则}：逐步增加 $r$ 直到感兴趣的响应量收敛。
\end{enumerate}

\subsection{模态参与因子与有效质量}

\textbf{模态参与因子}（Modal Participation Factor）衡量各阶模态对某一方向整体响应的贡献：

若基础运动方向为 $\mathbf{r}$（单位方向向量），第 $i$ 阶的模态参与因子为：

\begin{myformula}
\[
\Gamma_i = \frac{\bm{\phi}_i^{\top} \mathbf{M} \mathbf{r}}{m_i}
\]
\end{myformula}

\textbf{有效模态质量}（Effective Modal Mass）为：

\begin{myformula}
\[
m_{\text{eff},i} = \frac{(\bm{\phi}_i^{\top}\mathbf{M}\mathbf{r})^2}{m_i}
\]
\end{myformula}

重要性质：所有模态的有效质量之和等于刚体运动方向上的总质量：

\[
\sum_{i=1}^{n} m_{\text{eff},i} = \mathbf{r}^{\top}\mathbf{M}\mathbf{r} = M_{\text{total}}
\]

\textbf{工程截断标准}：使前 $r$ 阶的有效质量之和 $\geq 90\%$ 的总质量。这是结构抗震设计中模态截断的最常用准则。

\section{实验模态分析（EMA）简介}

\subsection{频响函数矩阵}

实验模态分析（Experimental Modal Analysis, EMA）是模态分析的实验对应。其核心是测量频响函数矩阵 $\mathbf{H}(\omega)$。

\textbf{定义}：$\mathbf{H}(\omega)$ 的 $(j,k)$ 元素 $H_{jk}(\omega)$ 表示仅在 $k$ 自由度上施加简谐力时，$j$ 自由度的稳态响应与力的复数比值。

\[
\hat{\mathbf{X}}(\omega) = \mathbf{H}(\omega)\,\hat{\mathbf{F}}(\omega)
\]

在模态坐标下，$\mathbf{H}(\omega)$ 可展开为各阶模态贡献的叠加：

\begin{myformula}
\[
\mathbf{H}(\omega) = \sum_{i=1}^{n} \frac{\bm{\phi}_i \bm{\phi}_i^{\top}}{-\omega^2 m_i + i\omega c_i + k_i} = \sum_{i=1}^{n} \frac{\bm{\phi}_i \bm{\phi}_i^{\top}}{m_i(\omega_i^2 - \omega^2 + 2i\zeta_i\omega_i\omega)}
\]
\end{myformula}

这一展开式即为 EMA 的数学基础。对高阻尼或密集模态，需要更精细的 MDOF 拟合方法。

\subsection{模态参数识别的三种方法}

\begin{enumerate}
\item \textbf{峰值拾取法（Peak Picking）}：最基本。FRF 幅值曲线上的峰 $\to \omega_i$，半功率带宽 $\to \zeta_i$，FRF 各通道在 $\omega_i$ 处的虚部之比 $\to$ 振型。适用于阻尼小、模态稀疏的情况。

\item \textbf{圆拟合（Circle Fit / Nyquist）}：单自由度拟合。在 Nyquist 图（实部 vs 虚部）中，模态响应近似为圆，圆的大小 $\to$ 阻尼，圆心角变化率最大处 $\to$ 固有频率。

\item \textbf{多自由度曲线拟合}：如 PolyMAX（Poly-reference Least-squares Complex Frequency-domain）、随机子空间法（SSI）等，可同时处理多模态、宽频带的 FRF 数据。
\end{enumerate}

\subsection{EMA 基本流程}

\begin{enumerate}
\item \textbf{激励}：力锤或激振器（选择合适的激励信号：随机、猝发随机、正弦扫频等）。
\item \textbf{测量}：加速度计/力传感器获取 FRF 或 IRF。
\item \textbf{信号处理}：$H_1$ / $H_2$ / $H_v$ 估计器降低噪声，相干函数检查数据质量。
\item \textbf{参数识别}：提取 $\omega_i$, $\zeta_i$, $\bm{\phi}_i$。
\item \textbf{验证}：MAC（模态置信准则）检查振型正交性和一致性，合成 FRF 与实测 FRF 比对。
\end{enumerate}

\section{Python 代码实现}

\subsection{$n$ 自由度系统模态分析}

\begin{mycode}{n自由度系统特征值求解与模态分析}
\begin{lstlisting}[language=Python]
import numpy as np
from scipy.linalg import eigh

def modal_analysis(M, K, n_modes=None):
    """
    n 自由度系统的模态分析
    
    Parameters
    ----------
    M, K : (n, n) ndarray  质量矩阵和刚度矩阵（对称）
    n_modes : int          需要保留的模态数（默认全部）
    
    Returns
    -------
    omega_n : ndarray      固有频率 (rad/s)
    phi : ndarray          模态矩阵 (质量归一化)
    f_n : ndarray          固有频率 (Hz)
    """
    n = M.shape[0]
    if n_modes is None:
        n_modes = n
    
    # 求解广义特征值问题 (仅前 n_modes 个)
    # eigh 返回从小到大排序的特征值和对应振型
    eigvals, eigvecs = eigh(K, M, subset_by_index=(0, n_modes - 1))
    
    # 固有频率
    omega_sq = np.maximum(eigvals, 0.0)  # 防止数值误差导致的微小负数
    omega_n = np.sqrt(omega_sq)
    f_n = omega_n / (2 * np.pi)
    
    # 振型质量归一化：phi^T * M * phi = I
    phi = np.zeros((n, n_modes))
    for i in range(n_modes):
        vec = eigvecs[:, i]
        norm_factor = np.sqrt(vec @ M @ vec)
        if norm_factor > 1e-15:
            phi[:, i] = vec / norm_factor
        else:
            phi[:, i] = vec
    
    return omega_n, phi, f_n


# ============ 示例：5 自由度弹簧-质量链 ============
n_dof = 5
m_val = 1.0   # 各质量相等
k_val = 1000.0 # 各弹簧相等

M = m_val * np.eye(n_dof)
K = np.zeros((n_dof, n_dof))
for i in range(n_dof):
    K[i, i] = 2 * k_val if 0 < i < n_dof - 1 else k_val
    if i > 0:
        K[i, i-1] = -k_val
    if i < n_dof - 1:
        K[i, i+1] = -k_val

omega_n, phi, f_n = modal_analysis(M, K)

print("固有频率:")
for i, (w, f) in enumerate(zip(omega_n, f_n)):
    print(f"  Mode {i+1}: omega = {w:8.4f} rad/s, f = {f:8.4f} Hz")

# 验证正交性
I_check = phi.T @ M @ phi
print(f"\n正交性: max|Phi^T M Phi - I| = {np.max(np.abs(I_check - np.eye(n_dof))):.2e}")

# 模态有效质量 (r = [1, 1, ..., 1]^T 方向)
r_dir = np.ones(n_dof)
for i in range(n_dof):
    gamma = phi[:, i] @ M @ r_dir
    m_eff = gamma ** 2
    print(f"  Mode {i+1}: m_eff = {m_eff:.4f}, "
          f"累计 = {sum((phi[:, :i+1].T @ M @ r_dir) ** 2):.4f}")
\end{lstlisting}
\end{mycode}

\subsection{模态叠加法时域响应（含 Rayleigh 阻尼）}

\begin{mycode}{n自由度系统模态叠加法时域求解}
\begin{lstlisting}[language=Python]
def modal_superposition_ndof(M, K, F, t, zeta, alpha=None, beta=None):
    """
    模态叠加法求解 n 自由度系统（Rayleigh 阻尼）
    
    F: (n_step, n_dof) 或 (n_step,) 外力数组
    zeta: 各阶模态阻尼比（如不详，用 alpha, beta 推算）
    alpha, beta: Rayleigh 阻尼系数
    """
    n_dof = M.shape[0]
    n_step = len(t)
    dt = t[1] - t[0]
    
    # 模态分析
    omega_n, phi, _ = modal_analysis(M, K)
    
    # 若未指定 zeta，由 Rayleigh 系数计算
    if zeta is None and alpha is not None and beta is not None:
        zeta = 0.5 * (alpha / omega_n + beta * omega_n)
    elif zeta is None:
        zeta = np.full(n_dof, 0.01)  # 默认 1%
    
    # 模态力
    f_modal = F @ phi  # (n_step, n_dof)
    
    # 各模态 Duhamel 积分数值解
    q = np.zeros((n_step, n_dof))
    q_dot = np.zeros((n_step, n_dof))
    
    for mode in range(n_dof):
        w = omega_n[mode]
        z = zeta[mode]
        wd = w * np.sqrt(1 - z**2) if z < 1 else 0.0
        f_m = f_modal[:, mode]
        
        if z >= 1:
            continue  # 临界/过阻尼另行处理
        
        for i in range(1, n_step):
            tau = t[i] - t[:i+1]
            h = np.exp(-z * w * tau) * np.sin(wd * tau) / wd
            if i > 0:
                q[i, mode] = np.trapz(f_m[:i+1] * h[::-1], dx=dt)
    
    # 变换回物理坐标
    x = (phi @ q.T).T  # (n_step, n_dof)
    return x, q, phi, omega_n


# 示例：5 自由度系统受半正弦脉冲
t_ex = np.arange(0, 3.0, 0.005)
F_ex = np.zeros((len(t_ex), n_dof))
td = 0.2
mask = (t_ex >= 0) & (t_ex <= td)
F_ex[mask, 0] = 1000.0 * np.sin(np.pi * t_ex[mask] / td)

x_modal, q_modal, phi_mat, omega_vec = modal_superposition_ndof(
    M, K, F_ex, t_ex, zeta=np.full(n_dof, 0.02))
\end{lstlisting}
\end{mycode}

\subsection{频响函数矩阵计算}

\begin{mycode}{频响函数矩阵计算}
\begin{lstlisting}[language=Python]
def compute_frf_matrix(omega_n, phi, zeta, omega_range):
    """
    由模态参数计算频响函数矩阵 H(omega)
    
    omega_n: (n_modes,)  固有频率
    phi: (n_dof, n_modes) 模态矩阵 (质量归一化)
    zeta: (n_modes,)     模态阻尼比
    omega_range: (n_freq,) 频率范围
    """
    n_dof, n_modes = phi.shape
    n_freq = len(omega_range)
    H = np.zeros((n_freq, n_dof, n_dof), dtype=complex)
    
    for i, w in enumerate(omega_range):
        for mode in range(n_modes):
            wm = omega_n[mode]
            zm = zeta[mode]
            denominator = wm**2 - w**2 + 2j * zm * wm * w
            H[i] += np.outer(phi[:, mode], phi[:, mode]) / denominator
    
    return H


# 计算示例：第 1 自由度激励 -> 第 3 自由度响应的 FRF
omega_scan = np.linspace(1, 200, 2000)
zeta_all = np.full(n_dof, 0.02)
H_mat = compute_frf_matrix(omega_n, phi, zeta_all, omega_scan)
FRF_1_3 = np.abs(H_mat[:, 2, 0])  # H_{3,1}(omega)
\end{lstlisting}
\end{mycode}

\section{小结}

本章将振动理论从两自由度推广至 $n$ 自由度，建立了完整的模态分析框架：

\begin{enumerate}
\item \textbf{运动方程}：$\mathbf{M}\ddot{\mathbf{x}} + \mathbf{C}\dot{\mathbf{x}} + \mathbf{K}\mathbf{x} = \mathbf{F}(t)$，矩阵形式统一描述了任意多自由度系统。

\item \textbf{模态分析}：通过求解广义特征值问题 $(\mathbf{K} - \omega^2\mathbf{M})\bm{\phi} = \mathbf{0}$，得到 $n$ 个固有频率和振型。

\item \textbf{正交性}：振型关于 $\mathbf{M}$ 和 $\mathbf{K}$ 正交，这是模态叠加法的理论基石。

\item \textbf{模态叠加法}：通过坐标变换 $\mathbf{x} = \mathbf{\Phi}\mathbf{q}$ 将耦合系统解耦为 $n$ 个独立的单自由度方程。

\item \textbf{比例阻尼}：$\mathbf{C} = \alpha\mathbf{M} + \beta\mathbf{K}$ 使模态叠加法得以在阻尼系统中直接应用。

\item \textbf{模态截断}：只取前 $r$ 阶模态近似系统响应，大幅降低计算成本。有效模态质量是指导截断的关键指标。

\item \textbf{实验模态分析}：频响函数的模态展开式为实验识别模态参数提供了理论工具。
\end{enumerate}

这一框架是所有现代结构动力学分析的基础，也是有限元动力分析和振动控制设计的核心工具。
